\section{模型结构}

自有模型分为锂盐和民爆两类利润池。锂盐收入由销量与实现价格构成，盈利能力由原料、加工与库存相关成本共同决定；民爆产品、爆破矿服和运输按历史收入、毛利与季节性建立保守基线。两类利润池汇总后，经费用、税费和股本口径形成归母与 EPS。正式 H1 预告只用来校准 2026 年上半年的合并利润区间，不作为销量、成本或全年利润的替代值。

\begin{keyinsight}[复算纪律]
所有读者可见的目标价均可由“情景 EPS $\times$ 周期校准估值倍数”及其市场/卖方锚权重复算。市场锚反映可观察的交易定价，外部卖方锚只采用完整原始报告；两者均不能取代本报告的基本面判断。
\end{keyinsight}

\section{情景与信心}

熊、基准和牛情景分别对应价格/销量/成本兑现较弱、符合审慎预期和行业条件显著改善。情景概率不是客观频率，而是截至数据截止日对可得证据和未披露字段的纪律化表达；其总和为 100\%。本报告不以资源、客户或项目的未披露增量来提高基准情景。具体假设、EPS、倍数、目标、概率及敏感性见第八章和结构化估值数据室。
